\itemhead{复习题三 · 第20题（教材p.92）}{XJ-B1-C03-RV-P092-Y1491}
\label{XJ-B1-C03-RV-P092-Y1491}
{\footnotesize 湘教版教材p.92 / PDF第99页 / 复习题三 / 第20题}
\field{原题}{探索函数 $y=x+\dfrac cx$（常数 $c>0$）的奇偶性、值域以及单调性，并说明理由；若函数为 $y=x^2+\dfrac c{x^2}$（常数 $c>0$），该函数的性质有何变化？}
\field{参考答案}{第一函数定义域 $\mathbb R\setminus\{0\}$，为奇函数，值域 $(-\infty,-2\sqrt c]\cup[2\sqrt c,+\infty)$；增区间 $(-\infty,-\sqrt c],[\sqrt c,+\infty)$，减区间 $[-\sqrt c,0),(0,\sqrt c]$。差式为 $(x_2-x_1)(1-c/(x_1x_2))$。第二函数为偶函数，值域 $[2\sqrt c,+\infty)$；令 $r=\sqrt[4]c$，增区间 $[-r,0),[r,+\infty)$，减区间 $(-\infty,-r],(0,r]$；用 $t=x^2>0$ 及基本不等式证明。两个断开的区间须分别写，不合并成单调区间。}
{\small \field{关联知识}{F08 奇偶性、对称与函数性质应用（教材帮p.113、114、115 / PDF122、123、124）；F05 单调性定义、证明与参数（教材帮p.97、98、100 / PDF106、107、109）；F06 最大值、最小值与值域（教材帮p.99、83、84 / PDF108、92、93）}
\field{方法与易错点}{检查定义域，再代入 $-x$；求另一侧解析式时同时写出 $-x$ 所在区间。 易错：只看公式不看定义域；用一对点证明奇偶；负半轴代换遗漏最外层负号。}
\field{答案依据}{答案册相关题组：PDF第24、25、26、27页（册内页23、24、25、26）；按当前题库原题逐项重算，不按旧版题号机械对套。}}
